\scanpage{54}

c) \deleted{\unsure{Soient}} $B$ un autre faisceau abélien. Supposons comme dans a) que \deleted{\unsure{le}} $\underline{\operatorname{Hom}}_{\mathrm{gr}}(A,C)=0$. Alors on trouve que la \deleted{\unsure{fonction}} \inserted{catégorie} \deleted{\unsure{est nulle}} $\underline{\operatorname{BIEXT}}(A,B;C)$ est discrète et
\[
\deleted{\operatorname{Biext}(A,B;C)\longleftarrow}
\]
\[
\operatorname{Biext}(A,B;C)\xrightarrow{\sim}\operatorname{Hom}_{\mathrm{gr}}\bigl(B,\overset{\inserted{A'}}{\underline{\operatorname{Ext}}^1}(A,C)\bigr).
\]
\inserted{\[
\operatorname{Biext}(A,B;C)\simeq H^0(S,\underline{\operatorname{Biext}}(A,B;C))\simeq\underline{\operatorname{Hom}}_{\mathrm{gr}}(B,\underline{\operatorname{Ext}}^1(A,C)).
\]}

Si on suppose \inserted{\unsure{de}} $\simeq$ que \deleted{\unsure{Biext}} $\underline{\operatorname{Hom}}_{\unsure{\mathrm{fct}}}(A,C)=0$, on trouve
\diagram{\begin{tikzcd}[column sep=small,row sep=large]
 & \operatorname{Hom}_{\unsure{\mathrm{ens}}}(B,R^1f_*(C_A)) \arrow[dl] \arrow[d,equal] & \\
\operatorname{Biext}(A,B;C) \arrow[r,leftrightarrow] & \operatorname{Hom}_{\mathrm{gr}}(B,R^1f_*(C_A)) \arrow[r] & H^0(B,R^1f_{B*}(C_{A\times_S B})) \arrow[d,equal] \\
 & \deleted{\unsure{f_*(C_A)}} & \begin{gathered}\text{classes de }C_{A\times_S B}\text{-torseurs}\\\deleted{\unsure{foncteur}}\\\deleted{\unsure{foncteurs}}\\\text{rigidifiés par}\ \unsure{\text{rapport à}}\ e_A\times_S B\end{gathered}\arrow[d,hook] \\
 & & H^{\unsure{1}}(A\times_S B,C_{A\times_S B})
\end{tikzcd}}
\margin{$=\text{classes de }C_{A\times_S B}\text{-torseurs rigidifiables par }e_A\times_S B$}

\diagram{\begin{tikzcd}
A \arrow[r,dash] \arrow[d,dash,"f"'] & A\times_S B \arrow[d,dash,"f_B"] \\
S \arrow[r,dash] & B
\end{tikzcd}}

Les torseurs sur $A\times_S B$ \deleted{\unsure{qui}} de gr $C_{A\times_S B}$ qu'on obtient sont ceux qui \inserted{a)} sont $e_A\times_S B$ rigidifiables; et pour lesquels l'application $B\mathbin{-}R^1f_*(C_A)$ qui le définit a deux vertus: b) c'est additif c) ça se factorise par
\[
\deleted{\unsure{R^1f_*}}\quad\underline{\operatorname{Ext}}^1(A,C)\simeq A'.
\]
\margin{\uline{Prop.} Les conditions a) et c) ensemble signifient que l'hom $A\mathbin{-}R^1q_*(C_B)$ est additif}

\deleted{Supposons qu'on ait également} \inserted{Notons d'ailleurs que la condition a) implique} que $L$ est \uline{birigidifiable} par. à $e_A,e_B$, et on trouve des foncteurs \uline{pl. fid.} de catégories discrètes
\diagram{\begin{tikzcd}[column sep=small]
\operatorname{BIEXT}(A,B;C) \arrow[r] & \operatorname{TORS}\,\unsure{\operatorname{RIG}}_{e_A,e_B}(A,B;C) \arrow[r] & \operatorname{TORS}\,\operatorname{RIG}_{e_A\times_S B}(C_{A\times_S B})
\end{tikzcd}}

La description \inserted{en termes du} \deleted{\unsure{par C torseurs de}} milieu est parfois préférable, car a) il est \deleted{\unsure{par}} symétrique en $A,B$, b) dans des cas importants, la pr. flèche est $\simeq$ une équivalence de catégorie.

\margin{53}
